% ============================================================
%  Symmetric-key encryption, perfect secrecy, One-Time Pad, Shannon
% ============================================================

\section{Symmetric-key encryption}\label{sec:ske}

Let's fix the scene that we will use for the whole chapter. \textbf{Alice} wants to send a message to \textbf{Bob} over a public channel. \textbf{Eve}, the eavesdropper, reads everything that travels on that channel. Alice and Bob have one advantage over Eve: they share a \textbf{secret key}, chosen uniformly at random, that Eve does not know. Since the \textbf{same key is used both to encrypt and to decrypt}, this setting is called \textbf{symmetric} (or private-key).

\begin{figure}[H]
  \centering
  \begin{tikzpicture}[node distance=1.2cm, box/.style={draw, rounded corners=2pt, minimum height=7mm, inner sep=4pt}]
    \node (A) {Alice};
    \node[box, right=1.2cm of A] (E) {$\Enc(k, m)$};
    \coordinate[right=1.6cm of E] (mid);
    \node[box, right=1.6cm of mid] (D) {$\Dec(k, c)$};
    \node[right=1.2cm of D] (B) {Bob};
    \node[below=1.1cm of mid] (Eve) {Eve};
    \node[below=0.9cm of E] (k1) {$k$};
    \node[below=0.9cm of D] (k2) {$k$};
    \draw[msg] (A) -- node[above] {$m$} (E);
    \draw[thick] (E) -- node[above] {$c$} (mid);
    \draw[msg] (mid) -- (D);
    \draw[msg, dashed] (mid) -- (Eve);
    \draw[msg] (D) -- node[above] {$m$} (B);
    \draw[msg] (k1) -- (E);
    \draw[msg] (k2) -- (D);
  \end{tikzpicture}
  \caption{Symmetric-key encryption. Eve sees the ciphertext $c$ but not the shared key $k$.}
\end{figure}

The objects we work with:
\begin{itemize}
  \item \textbf{Key space} $\Kspace$ $\to$ the set of possible keys; the key is sampled uniformly, $K \sample \Kspace$;
  \item \textbf{Message space} $\Mspace$ $\to$ the set of possible plaintexts;
  \item \textbf{Ciphertext space} $\Cspace$ $\to$ the set of possible ciphertexts.
\end{itemize}

\begin{frameddefn}[label={def:ske}]{Symmetric-key encryption scheme}
  A symmetric-key encryption scheme is a pair $\Pi = (\Enc, \Dec)$ with
  \[
    \Enc : \Kspace \times \Mspace \to \Cspace \qquad \Dec : \Kspace \times \Cspace \to \Mspace
  \]
  that satisfies \textbf{correctness}: decrypting with the right key always gives back the original message,
  \[
    \forall m \in \Mspace,\ \forall k \in \Kspace : \quad \Dec(k, \Enc(k, m)) = m
  \]
\end{frameddefn}

Correctness is only the minimum requirement: it says that the scheme \textit{works}, not that it is \textit{secure}. The scheme $\Enc(k, m) = m$ is perfectly correct and completely useless.

\subsection{Kerckhoffs's principle}

Where should the security of a scheme come from? Kerckhoffs's answer (1883), still valid today:

\begin{framedobs}{Kerckhoffs's principle}
  The security of a scheme must depend \textbf{only on the secrecy of the key}, never on the secrecy of the algorithm. Shannon summarized it as \textit{``the enemy knows the system''}.
\end{framedobs}

Why is this the right choice?
\begin{itemize}
  \item \textbf{Algorithms leak} $\to$ software gets reverse-engineered, devices get stolen, employees change jobs. A key is short and easy to replace; an algorithm is not.
  \item \textbf{Public algorithms get tested} $\to$ a public algorithm can be analysed by the whole research community, which is the best way to find weaknesses before an attacker does.
\end{itemize}
From now on, Eve always knows exactly how $\Enc$ and $\Dec$ work. The only thing she does not know is the key.

\section{Perfect secrecy}\label{sec:perfect-secrecy}

We now want to formalize the idea that ``Eve learns nothing from the ciphertext''. Think about what Eve knows \textbf{before} seeing anything: she is not completely ignorant. She may know that the message is probably ``attack at dawn'' and unlikely to be ``attack at noon''. This is her \textbf{a priori} knowledge, and we model it as a probability distribution over the messages.

Then Eve sees the ciphertext $c$ and updates her knowledge: this is the \textbf{a posteriori} knowledge. Shannon's idea (1949) is simple: \textbf{the scheme is perfectly secret if seeing the ciphertext does not change Eve's knowledge at all.}

The formal setup:
\begin{itemize}
  \item \textbf{Message} $\to$ a random variable $M$ over $\Mspace$ with an \textbf{arbitrary} distribution (what Eve knows a priori);
  \item \textbf{Key} $\to$ $K \sample \Kspace$, uniform and \textbf{independent of} $M$ (the key is chosen regardless of the message);
  \item \textbf{Ciphertext} $\to$ the random variable $C = \Enc(K, M)$ that we obtain from them.
\end{itemize}

\begin{frameddefn}[label={def:perfect-secrecy}]{Perfect secrecy}
  $\Pi$ is \textbf{perfectly secret} if, \textbf{for every distribution of $M$}, for every $m \in \Mspace$ and every $c \in \Cspace$ with $\Pr[C = c] > 0$:
  \[
    \Pr[M = m \mid C = c] = \Pr[M = m]
  \]
  In words: the a posteriori probability of every message is equal to its a priori probability.
\end{frameddefn}

Two details of this definition are essential:
\begin{itemize}
  \item \textbf{``For every distribution of $M$''} $\to$ security is a property of the \textbf{scheme}, and it must hold whatever Eve already knows about the message. We will see in \Cref{sec:equivalent} that without this quantifier the definition would be meaningless.
  \item \textbf{The adversary is unbounded} $\to$ nothing in the definition talks about time or resources: Eve can compute anything she wants.
\end{itemize}

\paragraph{Notation used in class.} The joint probability is written with the symbol $\wedge$ (``and''): $\Pr[M = m \wedge C = c]$ is the probability that the message is $m$ \textbf{and} the ciphertext is $c$. It is the same as $\Pr[M = m, C = c]$. With this notation, the left side of the definition is computed as
\[
  \Pr[M = m \mid C = c] = \frac{\Pr[M = m \wedge C = c]}{\Pr[C = c]}
\]
Knowing that the ciphertext is $c$ \textbf{restricts the world} to the outcomes where the ciphertext is $c$, and among those we measure in which fraction the message is $m$.

\begin{framedex}{Perfect secrecy with numbers}
  A general sends one bit with the One-Time Pad: ``attack'' (0) with probability $0.9$, ``retreat'' (1) with probability $0.1$. The key is 0 or 1, each with probability $0.5$. Eve sees the ciphertext 0. How likely is ``attack''?
  \begin{itemize}
    \item \textbf{Numerator} $\to$ ``attack'' \textbf{and} ciphertext 0 needs key 0: $\Pr[M = \text{attack} \wedge C = 0] = 0.9 \cdot 0.5 = 0.45$;
    \item \textbf{Denominator} $\to$ the ciphertext is 0 in two cases, ``attack'' with key 0 ($0.45$) or ``retreat'' with key 1 ($0.1 \cdot 0.5 = 0.05$), so $\Pr[C = 0] = 0.5$.
  \end{itemize}
  \[
    \Pr[M = \text{attack} \mid C = 0] = \frac{0.45}{0.5} = 0.9
  \]
  Before seeing the ciphertext Eve believed ``attack'' with probability $0.9$, and after seeing it she still believes $0.9$: \textbf{she learned nothing}. With a scheme that does not encrypt at all (ciphertext equal to the message) she would get $\frac{0.9}{0.9} = 1$: she would learn everything.
\end{framedex}

\section{Three equivalent definitions of perfect secrecy}\label{sec:equivalent}

The definition above describes what we want, but it is not very handy when we have to check a concrete scheme. The following lemma gives two equivalent ways to say the same thing, and the third one will be our main tool.

\begin{framedlem}[label={lem:equivalent}]{Equivalent definitions of perfect secrecy}
  For every symmetric-key encryption scheme $\Pi$, the following statements are equivalent:
  \begin{enumerate}[label=(\roman*)]
    \item \textbf{Perfect secrecy}: $\Pr[M = m \mid C = c] = \Pr[M = m]$ for every distribution of $M$;
    \item \textbf{Independence}: $M$ and $C$ are independent for every distribution of $M$, that is,
    \[
      \Pr[M = m \wedge C = c] = \Pr[M = m] \cdot \Pr[C = c]
    \]
    \item \textbf{Same ciphertext distribution}: for all $m, m' \in \Mspace$ and all $c \in \Cspace$,
    \[
      \Pr_K[\Enc(K, m) = c] = \Pr_K[\Enc(K, m') = c]
    \]
  \end{enumerate}
\end{framedlem}

How to read the three statements:

\begin{table}[H]
  \centering
  \small
  \begin{tabularx}{\textwidth}{@{}lXX@{}}
    \toprule
    \textbf{Statement} & \textbf{What it talks about} & \textbf{When it is useful} \\
    \midrule
    (i) Perfect secrecy & Eve's knowledge, a priori versus a posteriori & the \textit{intuitive} definition: what we actually want \\
    (ii) Independence & the joint distribution of $M$ and $C$ & a bridge between (i) and (iii); in information theory it means that the mutual information is zero, $I(M;C) = 0$ \\
    (iii) Same ciphertext distribution & only the randomness of the \textbf{key}: no distribution of $M$ is involved & the \textit{operational} one: the easiest to check on a concrete scheme \\
    \bottomrule
  \end{tabularx}
\end{table}

Statement (iii) deserves one more sentence: \textbf{the distribution of the ciphertext does not depend on the message}. Whatever message Alice encrypts, the ciphertext ``looks the same'' to Eve, so observing it cannot tell her which message was sent.

We prove the equivalence with a cycle of implications: (i) $\Rightarrow$ (ii) $\Rightarrow$ (iii) $\Rightarrow$ (i).

\begin{proof}[Proof of (i) $\Rightarrow$ (ii)]
  \textbf{Goal:} starting from perfect secrecy, obtain $\Pr[M = m \wedge C = c] = \Pr[M = m] \cdot \Pr[C = c]$.

  We start from hypothesis (i) and expand the conditional probability with its definition:
  \[
    \Pr[M = m] \overset{\text{(i)}}{=} \Pr[M = m \mid C = c] = \frac{\Pr[M = m \wedge C = c]}{\Pr[C = c]}
  \]
  The first equality is hypothesis (i); the second is the definition of conditional probability, which can be used because $\Pr[C = c] > 0$. Multiplying both sides by $\Pr[C = c]$:
  \[
    \Pr[M = m \wedge C = c] = \Pr[M = m] \cdot \Pr[C = c]
  \]
  which is exactly independence. If $\Pr[C = c] = 0$, both sides are $0$ and the equality holds anyway.
\end{proof}

\begin{proof}[Proof of (ii) $\Rightarrow$ (iii)]
  \textbf{Goal:} show that $\Pr[\Enc(K, m) = c] = \Pr[\Enc(K, m') = c]$ for any two messages. \textbf{Idea:} show that \textbf{both} are equal to the same quantity, $\Pr[C = c]$. We start from the left side and transform it step by step:
  \[
    \Pr[\Enc(K, m) = c] \overset{(a)}{=} \Pr[\Enc(K, M) = c \mid M = m] \overset{(b)}{=} \Pr[C = c \mid M = m] \overset{(c)}{=} \Pr[C = c]
  \]
  \begin{itemize}
    \item \textbf{(a)} ``Encrypting the \textit{fixed} message $m$'' is the same as ``encrypting the \textit{random} message $M$, knowing that it turned out to be $m$''. This works because the key $K$ is independent of $M$: knowing that $M = m$ tells us nothing about the key.
    \item \textbf{(b)} By definition $C = \Enc(K, M)$, so we are just renaming.
    \item \textbf{(c)} Hypothesis (ii): $M$ and $C$ are independent, so knowing that $M = m$ does not change the probability that $C = c$.
  \end{itemize}
  By the same steps with $m'$, $\Pr[\Enc(K, m') = c] = \Pr[C = c]$. Both probabilities are equal to $\Pr[C = c]$, so they are equal to each other.

  A technical detail: to condition on $M = m$ we need $\Pr[M = m] > 0$. Since (ii) holds \textbf{for every distribution}, we can choose $M$ uniform, so that every message has positive probability.
\end{proof}

\begin{proof}[Proof of (iii) $\Rightarrow$ (i)]
  \textbf{Goal:} from ``the ciphertext has the same distribution for every message'', obtain perfect secrecy. Fix any distribution of $M$, and $m, c$ with $\Pr[C = c] > 0$. The proof has two steps:
  \begin{enumerate}
    \item show that $\Pr[C = c] = \Pr[C = c \mid M = m]$: knowing the message does not change the probability of the ciphertext;
    \item use Bayes' rule to ``reverse'' this equality and obtain $\Pr[M = m] = \Pr[M = m \mid C = c]$.
  \end{enumerate}

  \textbf{Step 1.} We start from $\Pr[C = c]$ and transform it step by step:
  \[
    \begin{aligned}
      \Pr[C = c] &\overset{(a)}{=} \sum_{m'} \Pr[C = c \wedge M = m'] \\
      &\overset{(b)}{=} \sum_{m'} \Pr[C = c \mid M = m'] \cdot \Pr[M = m'] \\
      &\overset{(c)}{=} \sum_{m'} \Pr[\Enc(K, M) = c \mid M = m'] \cdot \Pr[M = m'] \\
      &\overset{(d)}{=} \sum_{m'} \Pr[\Enc(K, m') = c] \cdot \Pr[M = m'] \\
      &\overset{(e)}{=} \sum_{m'} \Pr[\Enc(K, m) = c] \cdot \Pr[M = m'] \\
      &\overset{(f)}{=} \Pr[\Enc(K, m) = c] \cdot \underbrace{\sum_{m'} \Pr[M = m']}_{=\,1} \\
      &\overset{(g)}{=} \Pr[\Enc(K, M) = c \mid M = m] = \Pr[C = c \mid M = m]
    \end{aligned}
  \]
  \begin{itemize}
    \item \textbf{(a)} Law of total probability: we split the event $C = c$ according to which message $m'$ was sent.
    \item \textbf{(b)} Definition of conditional probability, in the form $\Pr[A \wedge B] = \Pr[A \mid B] \cdot \Pr[B]$.
    \item \textbf{(c)} We write $C$ as what it is, $\Enc(K, M)$.
    \item \textbf{(d)} The same trick as step (a) of the previous proof: knowing that $M = m'$, encrypting $M$ is the same as encrypting the fixed message $m'$ (the key is independent of the message).
    \item \textbf{(e)} \textbf{Here we use hypothesis (iii)}: $\Pr[\Enc(K, m') = c]$ is the same for every message, so we can replace $m'$ with our fixed $m$.
    \item \textbf{(f)} $\Pr[\Enc(K, m) = c]$ no longer depends on $m'$, so it comes out of the sum. What is left is the sum of the probabilities of all messages, which is $1$.
    \item \textbf{(g)} The trick of (d) and (c) backwards: we go back to $\Pr[C = c \mid M = m]$.
  \end{itemize}

  \textbf{Step 2 (Bayes).} The probability of ``both events'' can be written in two ways, conditioning on one event or on the other:
  \[
    \Pr[M = m \mid C = c] \cdot \Pr[C = c] = \Pr[M = m \wedge C = c] = \Pr[C = c \mid M = m] \cdot \Pr[M = m]
  \]
  We solve for $\Pr[M = m]$ by dividing by $\Pr[C = c \mid M = m]$, which by Step 1 is equal to $\Pr[C = c] > 0$:
  \[
    \begin{aligned}
      \Pr[M = m] &= \frac{\Pr[M = m \mid C = c] \cdot \Pr[C = c]}{\Pr[C = c \mid M = m]} \\
      &= \frac{\Pr[M = m \mid C = c] \cdot \cancel{\Pr[C = c]}}{\cancel{\Pr[C = c]}} = \Pr[M = m \mid C = c]
    \end{aligned}
  \]
  This is perfect secrecy.
\end{proof}

\begin{framedobs}{Why ``for every distribution'' is essential}
  Suppose we dropped the quantifier and asked the definition to hold for \textbf{one single} distribution, for example the constant one $\Pr[M = m_0] = 1$ (Eve already knows the message). Then:
  \begin{itemize}
    \item statements (i) and (ii) would be true for \textbf{any} scheme, even $\Enc(k, m) = m$, which does not encrypt at all;
    \item statement (iii) would be false for that scheme.
  \end{itemize}
  So the lemma would be false. The step (ii) $\Rightarrow$ (iii) works only because we are free to choose the uniform distribution.
\end{framedobs}

\section{The One-Time Pad}\label{sec:otp}

We now have a definition and a handy way to check it. Does any scheme actually achieve perfect secrecy? Yes, and a very simple one: the \textbf{One-Time Pad} (Vernam, 1917). The idea is to hide every bit of the message by combining it with a \textbf{fresh, uniformly random key bit}, using the XOR operation (bitwise exclusive OR: $0$ if the two bits are equal, $1$ if they are different).

\begin{frameddefn}[label={def:otp}]{One-Time Pad}
  $\Mspace = \Kspace = \Cspace = \{0,1\}^n$, with the key sampled uniformly, $K \sample \{0,1\}^n$:
  \[
    \Enc(k, m) = k \oplus m \qquad \Dec(k, c) = k \oplus c
  \]
\end{frameddefn}

\textbf{Correctness.} We only need two properties of XOR: it is associative, and anything XOR-ed with itself gives zero ($k \oplus k = 0^n$):
\[
  \Dec(k, \Enc(k, m)) = k \oplus (k \oplus m) = (k \oplus k) \oplus m = 0^n \oplus m = m
\]

\begin{framedex}{One-Time Pad with four bits}
  Take $m = 1011$ and $k = 0110$. Then $c = 1101$, and decrypting gives back $c \oplus k = 1101 \oplus 0110 = 1011 = m$.

  Now look at it from Eve's side. She sees $c = 1101$. Could the message be $m' = 0000$? Yes: with the key $k' = m' \oplus c = 1101$. The same holds for \textbf{every} possible message: for each one there is \textbf{exactly one key} that produces $c$. Since all keys are equally likely, Eve has no reason to prefer one message over another. This is the intuition behind the proof.
\end{framedex}

\begin{framedprop}[label={prop:otp}]{}
  The One-Time Pad is perfectly secret.
\end{framedprop}

\begin{proof}
  We use statement (iii) of \Cref{lem:equivalent}, because it only involves the key. Fix any message $m$ and any ciphertext $c$:
  \[
    \Pr_K[\Enc(K, m) = c] \overset{(a)}{=} \Pr_K[K \oplus m = c] \overset{(b)}{=} \Pr_K[K = m \oplus c] \overset{(c)}{=} 2^{-n}
  \]
  \begin{itemize}
    \item \textbf{(a)} This is the definition of encryption in the One-Time Pad.
    \item \textbf{(b)} We ``solve for $K$'' by XOR-ing both sides with $m$: since $m \oplus m = 0^n$, we get $K \oplus m = c \iff K = m \oplus c$. So the event ``encrypting $m$ gives $c$'' is \textbf{exactly} the event ``the key happens to be $m \oplus c$'': there is one and only one key that does it.
    \item \textbf{(c)} $m \oplus c$ is one specific string and the key is uniform, so the probability of picking exactly that key is $\frac{1}{|\Kspace|} = 2^{-n}$.
  \end{itemize}
  The result \textbf{does not depend on $m$}: for any other message $m'$ we get $2^{-n}$ again. So statement (iii) holds and, by the lemma, the One-Time Pad is perfectly secret.
\end{proof}

\begin{framedex}{Intuition with one bit}
  With $n = 1$, ``attack'' $= 0$ and ``retreat'' $= 1$:
  \begin{center}
    \begin{tabular}{lcc}
      \toprule
      \textbf{Message} & \textbf{Key 0} & \textbf{Key 1} \\
      \midrule
      attack (0) & ciphertext 0 & ciphertext 1 \\
      retreat (1) & ciphertext 1 & ciphertext 0 \\
      \bottomrule
    \end{tabular}
  \end{center}
  If Eve sees ciphertext 0, it could be ``attack'' with key 0 or ``retreat'' with key 1, and both keys have probability $\frac{1}{2}$. The ciphertext gives her no reason to prefer one message over the other.
\end{framedex}

A useful byproduct: since $\Pr[C = c] = 2^{-n}$ for every $c$, the ciphertext is \textbf{uniformly random}, whatever the distribution of the message. To Eve it looks like pure noise.

\subsection{Limitations of the One-Time Pad}

If the One-Time Pad is perfect, why don't we use it everywhere? Because perfection has a price:

\begin{table}[H]
  \centering
  \small
  \begin{tabularx}{\textwidth}{@{}>{\raggedright\arraybackslash}p{2.9cm}>{\raggedright\arraybackslash}p{4.2cm}X@{}}
    \toprule
    \textbf{Limitation} & \textbf{Why} & \textbf{Consequence} \\
    \midrule
    Key as long as the message & $|\Kspace| = |\Mspace| = 2^n$ & To send $n$ bits, Alice and Bob must have \textit{already} shared $n$ secret random bits. If they can share the key securely, they could share the message the same way. \\
    One-time only & $c_1 \oplus c_2 = m_1 \oplus m_2$ & Reusing the key reveals the XOR of the two messages; with natural-language text both messages can often be recovered (historical case: the Venona project). \\
    Malleable & $\Dec(k, c \oplus \Delta) = m \oplus \Delta$ & Eve can flip chosen bits of the message without knowing the key: secrecy is not integrity, and this motivates MACs. \\
    Truly uniform key & the proof uses $\Pr[K = k] = 2^{-n}$ & With a biased key, perfect secrecy no longer holds. \\
    \bottomrule
  \end{tabularx}
\end{table}

\section{Shannon's theorem: perfect secrecy requires long keys}

The first limitation is the most serious one. Is it a defect of the One-Time Pad that a smarter scheme could fix? Shannon proved that the answer is \textbf{no}: it holds for \textbf{every} perfectly secret scheme.

\begin{framedthm}[label={thm:shannon}]{Shannon}
  If $\Pi = (\Enc, \Dec)$ is a correct symmetric-key encryption scheme with perfect secrecy, then
  \[
    |\Kspace| \geq |\Mspace|
  \]
  There must be at least as many keys as messages. In terms of lengths: if messages are $n$-bit strings, keys must have \textbf{at least} $n$ bits (equality is allowed, and the One-Time Pad achieves it).
\end{framedthm}

\begin{framedex}{The intuition with three messages and two keys}
  Suppose there are three messages (``attack'', ``retreat'', ``wait'') and only two keys, $k_1$ and $k_2$. Eve intercepts a ciphertext $c$ and, being unbounded, decrypts it with \textbf{every} key: with $k_1$ she gets, say, ``attack'' and with $k_2$ she gets ``retreat''. Since there are only two keys, there are at most two candidate messages, so ``wait'' is \textbf{excluded}: no key decrypts $c$ to ``wait'', so ``wait'' cannot be the message behind $c$.

  Eve has learned something (``it is not \textit{wait}''): before seeing $c$, ``wait'' had some positive probability; after seeing $c$, its probability is zero. Too few keys cannot ``cover'' all messages.
\end{framedex}

\begin{proof}
  By contradiction, suppose that $|\Kspace| < |\Mspace|$.
  \begin{enumerate}
    \item \textbf{Choose the distribution.} Perfect secrecy holds for every distribution, so we choose $M$ uniform over $\Mspace$: every message has probability $\frac{1}{|\Mspace|} > 0$.
    \item \textbf{Fix a ciphertext.} Take any $c$ with $\Pr[C = c] > 0$ (such a $c$ exists, since $C$ always takes some value).
    \item \textbf{Collect the possible decryptions.} Decrypting $c$ with every key gives the set
    \[
      \Mspace(c) := \{\, \Dec(k, c) : k \in \Kspace \,\} \qquad \text{with} \qquad |\Mspace(c)| \leq |\Kspace| < |\Mspace|
    \]
    since every key produces a single message.
    \item \textbf{Find a missing message.} $\Mspace(c)$ is smaller than $\Mspace$, so there is a message $m \in \Mspace \setminus \Mspace(c)$.
    \item \textbf{This message never produces $c$.} By correctness, if $\Enc(k, m) = c$ for some key $k$, then $\Dec(k, c) = m$, that is, $m \in \Mspace(c)$. This is impossible, so no key maps $m$ to $c$.
    \item \textbf{Contradiction.} Then
    \[
      \Pr[M = m \mid C = c] = 0 \neq \frac{1}{|\Mspace|} = \Pr[M = m]
    \]
    which contradicts perfect secrecy. \qedhere
  \end{enumerate}
\end{proof}

Consequences:
\begin{itemize}
  \item \textbf{The One-Time Pad is optimal} $\to$ it has $|\Kspace| = |\Mspace| = 2^n$, the smallest key space allowed by the theorem.
  \item \textbf{Perfect secrecy is inherently impractical} $\to$ to encrypt 1 gigabyte with perfect secrecy, Alice and Bob must have already shared 1 gigabyte of secret, uniformly random key, and they can never reuse it. To protect large amounts of data with short keys (for example 128 bits) we have to weaken the model: an adversary with \textbf{bounded computational power} and \textbf{hardness assumptions}. This is \textbf{computational security}.
\end{itemize}

\section{Going further}\label{sec:going-further}

\subsection{Shannon's characterization theorem}

When the three spaces have the same size, there is a very quick test for perfect secrecy (Katz and Lindell, Theorem 2.10).

\begin{framedthm}{Shannon's characterization}
  If $|\Mspace| = |\Kspace| = |\Cspace|$ and $\Enc$ is deterministic, then $\Pi$ is perfectly secret \textbf{if and only if}:
  \begin{enumerate}
    \item every key is chosen with probability exactly $1/|\Kspace|$;
    \item for every $m \in \Mspace$ and $c \in \Cspace$ there is a \textbf{unique} key $k$ such that $\Enc(k, m) = c$.
  \end{enumerate}
\end{framedthm}

The One-Time Pad passes both conditions immediately: keys are uniform, and $k = m \oplus c$ is the unique key.

\subsection{Exercise: a game-based definition of perfect secrecy}\label{sec:ps-game}

In class this was given as an exercise. The professor's formulation does not choose the bit $b$ at random: it \textbf{fixes} it, and compares two versions of the same game.

\begin{frameddefn}[label={def:ps-game}]{The games $\mathrm{GAME}_{\Pi,\Adv}(b)$}
  For a scheme $\Pi$, an adversary $\Adv$ (even unbounded) and a \textbf{fixed} bit $b \in \{0,1\}$:
  \begin{enumerate}
    \item $\Adv$ outputs two messages $m_0, m_1 \in \Mspace$;
    \item the key is sampled, $k \sample \Kspace$, and $\Adv$ receives $c = \Enc(k, m_b)$;
    \item $\Adv$ outputs a bit $b'$, which is the output of the game.
  \end{enumerate}
  $\Pi$ is \textbf{perfectly secret in the game sense} if for every adversary $\Adv$
  \[
    \Pr\big[\mathrm{GAME}_{\Pi,\Adv}(0) = 1\big] - \Pr\big[\mathrm{GAME}_{\Pi,\Adv}(1) = 1\big] = 0
  \]
  that is, the output of the game has \textbf{the same distribution} whether $m_0$ or $m_1$ is encrypted.
\end{frameddefn}

\begin{framedprob}[label={prob:ps-game}]{}
  Prove that $\Pi$ is perfectly secret (\Cref{def:perfect-secrecy}) if and only if it is perfectly secret in the game sense (\Cref{def:ps-game}).
\end{framedprob}

\begin{proof}[Solution]
  By \Cref{lem:equivalent}, perfect secrecy is equivalent to statement (iii): for all $m, m', c$, $\Pr_K[\Enc(K, m) = c] = \Pr_K[\Enc(K, m') = c]$. So it is enough to show that (iii) is equivalent to the game definition.

  \textbf{(iii) $\Rightarrow$ game.} Fix any adversary $\Adv$ and the pair $m_0, m_1$ it outputs. By (iii), the ciphertext $c = \Enc(K, m_0)$ and the ciphertext $c = \Enc(K, m_1)$ have \textbf{exactly the same distribution}. In both games, $\Adv$ computes $b'$ from $c$ (and from its own random bits, which do not depend on $b$). Applying the same procedure to the same distribution gives the same distribution, so $\Pr[b' = 1]$ is the same in the two games and the difference is $0$. If $\Adv$ chooses $m_0, m_1$ at random, the argument holds for every possible choice, and so also on average.

  \textbf{Game $\Rightarrow$ (iii).} By contradiction, suppose that (iii) fails: there are $m_0, m_1, c^*$ with
  \[
    \Pr_K[\Enc(K, m_0) = c^*] \neq \Pr_K[\Enc(K, m_1) = c^*]
  \]
  Consider the adversary $\Adv$ that outputs exactly these $m_0, m_1$, and outputs $b' = 1$ if and only if the ciphertext it receives is $c^*$. Then
  \[
    \Pr\big[\mathrm{GAME}(0) = 1\big] - \Pr\big[\mathrm{GAME}(1) = 1\big] = \Pr_K[\Enc(K, m_0) = c^*] - \Pr_K[\Enc(K, m_1) = c^*] \neq 0
  \]
  which contradicts the game definition. (This adversary is even efficient: it only compares $c$ with a fixed string. The unbounded power is needed only to \textit{find} $m_0, m_1, c^*$, and the definition quantifies over every adversary, so we may assume it knows them.)
\end{proof}

\paragraph{Link with perfect indistinguishability.} The game definition is also equivalent to perfect indistinguishability (\Cref{def:perfect-indist}), where $b$ is uniform and the adversary must guess it. Write $p_b = \Pr[\mathrm{GAME}(b) = 1]$. When $b$ is uniform:
\[
  \Pr[b' = b] = \frac{1}{2}\Pr[b' = 0 \mid b = 0] + \frac{1}{2}\Pr[b' = 1 \mid b = 1] = \frac{1}{2}(1 - p_0) + \frac{1}{2}p_1 = \frac{1}{2} + \frac{1}{2}(p_1 - p_0)
\]
So $\Pr[b' = b] = \frac{1}{2}$ for every adversary if and only if $p_0 = p_1$ for every adversary. The three formulations (perfect secrecy, the game with fixed $b$, the game with random $b$) say the same thing (see also Katz and Lindell, Lemma 2.6).

This is why the game-based formulation is so important: \textbf{all the computational definitions of the course follow the same pattern}, replacing ``every adversary'' with ``every probabilistic polynomial-time adversary'' and $\frac{1}{2}$ with $\frac{1}{2} + \negl(n)$.

\subsection{The entropy formulation}

The \textit{entropy} $H(X)$ measures how much uncertainty (randomness) a variable contains; the \textit{mutual information} $I(X;Y)$ measures how much information one variable gives about another. Statement (ii) of the lemma says exactly that the mutual information between message and ciphertext is zero, $I(M;C) = 0$. Shannon's theorem also has an \textbf{entropy version}, $H(K) \geq H(M)$:
\[
  H(M) = H(M \mid C) \leq H(M, K \mid C) = H(K \mid C) + \underbrace{H(M \mid K, C)}_{=\,0 \text{ (correctness)}} \leq H(K)
\]
It is stronger than $|\Kspace| \geq |\Mspace|$: it is not enough to have many possible keys, they must also contain at least as much \textit{randomness} as the message.

\subsection{Example and counterexample: the shift cipher}

\begin{framedex}{The shift cipher}
  $\Mspace = \Kspace = \Cspace = \Z_{26}$ (the letters of the alphabet) and $\Enc(k, m) = (m + k) \bmod 26$.
  \begin{itemize}
    \item \textbf{One letter} $\to$ $\Pr_K[\Enc(K, m) = c] = \Pr_K[K = c - m] = 1/26$ for all $m, c$ $\to$ \textbf{perfectly secret}. It is just the One-Time Pad over $\Z_{26}$.
    \item \textbf{Two letters with the same key} $\to$ the ciphertext ``XX'' can come from ``AA'' but \textbf{never} from ``AB'', because the shift preserves the distance between the letters. Statement (iii) fails $\to$ \textbf{not} perfectly secret. And indeed $|\Kspace| = 26 < 26^2 = |\Mspace|$, exactly as Shannon's theorem predicts.
  \end{itemize}
\end{framedex}
